\slajdid{09_1-s016}
\begin{frame}[label=09_1-s016]
Na osnovu prethodne analize možemo odmah da kažemo da je
\[V_{OH}=V_{DD}\]
dok ćemo napon $V_{OL}$ naći u trećoj oblasti zamenjujući $V_I=V_{OH}=V_{DD}$
\[\frac{V_{DD}-V_{OL}}{R_L}=\frac{k_n}{2}\frac1{1+\frac{V_{OL}}{L_nE_{Cn}}}\bigl(2V_{OL}(V_{DD}-V_{Tn})-V_{OL}^2\bigr)\]
I sada ide nešto što vas zbunjuje.\par\medskip
Možemo ovaj izraz „tačno“ rešiti. Kvadratna jednačina. \alert{Ali to nećemo raditi.}\par\medskip
\alert{Radićemo zanemarivanja}, koja će nas često dovesti do izraza za tranzistor sa dugačkim kanalom pa i dalje.\par\medskip
Prvo: nas interesuje kvalitativna slika uz kvantitativnu podršku - da malo grešimo.\par\medskip
Drugo: ovo ćemo zaista često raditi ali \alert{će postojati prilike kada to ne smemo da radimo}.
\end{frame}
